\documentclass[11pt,a4paper]{article}
\usepackage[margin=2.5cm]{geometry}
\usepackage{booktabs}
\usepackage{xcolor}
\usepackage{hyperref}
\usepackage{enumitem}
\usepackage{titlesec}
\usepackage{fancyhdr}
\usepackage{microtype}
\usepackage{parskip}
\usepackage{fontspec}
\usepackage{polyglossia}

\setmainlanguage{english}
\setotherlanguage{hebrew}
\newfontfamily\hebrewfont{DejaVu Sans}[Script=Hebrew]

\hypersetup{colorlinks=true, linkcolor=blue, urlcolor=blue}

\pagestyle{fancy}
\fancyhf{}
\rhead{Knesset Protocol Journal Project}
\lhead{Data \& Preprocessing Report}
\rfoot{Page \thepage}
\lfoot{Tomer Morad, Or Israeli, Yoel Marko --- June 2026}

\titleformat{\section}{\large\bfseries}{}{0em}{}[\titlerule]
\titleformat{\subsection}{\normalsize\bfseries}{}{0em}{}

\begin{document}

\begin{center}
  {\LARGE \textbf{Knesset Protocol Journal Project}}\\[0.4em]
  {\large Data Overview \& Preprocessing Report}\\[0.6em]
  {\normalsize Tomer Morad \quad Or Israeli \quad Yoel Marko}\\
  {\normalsize June 2026}
\end{center}

\vspace{1em}

% -----------------------------------------------------------------------
\section{Dataset Overview}

The raw dataset consists of \textbf{213 JSON files}, each representing a single Israeli
Knesset (25th) Finance Committee
(\texthebrew{ועדת הכספים}) meeting. Files follow the naming
convention \texttt{25\_ptv\_<doc\_id>.json}.

\subsection*{Per-file Schema}

\begin{center}
\begin{tabular}{lll}
\toprule
\textbf{Field} & \textbf{Type} & \textbf{Description} \\
\midrule
\texttt{knesset\_num} & int & Knesset session number (25) \\
\texttt{committee}    & string & Committee name (Hebrew) \\
\texttt{doc\_id}      & string & Unique protocol identifier \\
\texttt{date}         & ISO 8601 & Meeting date and time \\
\texttt{source\_file} & URL & Source document on Knesset website \\
\texttt{text}         & string & Full meeting transcript (Hebrew) \\
\bottomrule
\end{tabular}
\end{center}

\textbf{Language:} All transcript text is in Hebrew.
\textbf{Date range:} November 2022 -- June 2023.

% -----------------------------------------------------------------------
\section{Topic Header Extraction}

\subsection*{Primary Marker: \texttt{<{}< נושא >{}>}}

The vast majority of protocols embed topic headers using the pattern:

\begin{center}
\texttt{<{}< נושא >{}>} \emph{topic text} \texttt{<{}< נושא >{}>}
\end{center}

Each topic appears twice per file (once in the agenda block, once as a section header
in the transcript body). We extract unique spans, deduplicate within file, and strip
leading list numbering (e.g.\ \texttt{1. }, \texttt{2. }) from multi-topic agendas.

\subsection*{Fallback 1: \texttt{<{}< הצח >{}>} (Budget Law Variant)}

Four protocols use an alternate marker \texttt{<{}< הצח >{}>} instead of
\texttt{<{}< נושא >{}>}. These are sessions dedicated to the 2023--2024 State Budget Law,
which uses a different typographical convention in the official documents.
We apply identical extraction logic to this marker.

\subsection*{Fallback 2: \texthebrew{סדר היום} Agenda Block}

Four protocols carry no inline markers at all. For these, we extract the topic from the
\texthebrew{סדר היום} (``agenda'') block at the top of every transcript,
delimited between the meeting header and the \texthebrew{נכחו} (``attendees'') section.

\subsection*{Extraction Results}

\begin{center}
\begin{tabular}{lrr}
\toprule
\textbf{Method} & \textbf{Files} & \textbf{\%} \\
\midrule
\texttt{<{}< נושא >{}>} marker & 205 & 96.2\% \\
\texttt{<{}< הצח >{}>} marker  & 4   & 1.9\%  \\
\texthebrew{סדר היום} fallback & 4   & 1.9\%  \\
No topic found                 & 0   & 0.0\%  \\
\midrule
\textbf{Total}                 & \textbf{213} & \textbf{100\%} \\
\bottomrule
\end{tabular}
\end{center}

\textbf{All 213 protocols were successfully assigned at least one topic.}

% -----------------------------------------------------------------------
\section{Topic Canonicalization}

\subsection*{Raw vs.\ Canonical}

Raw extraction yielded \textbf{180 unique topic strings} across all 213 files.
Many are near-duplicates caused by:

\begin{itemize}[noitemsep]
  \item Typographical errors
    (\texthebrew{מיסים} vs.\ \texthebrew{מסים};
     \texthebrew{סגירה} vs.\ \texthebrew{סגירת};
     \texthebrew{להתקשרות} vs.\ \texthebrew{להתקשורת})
  \item Punctuation/dash variants (em-dash vs.\ regular hyphen)
  \item Bill number formatting
    (\texttt{(מ/1577)} vs.\ \texttt{, מ/1577} vs.\ \texttt{---2022, מ/1577})
  \item Missing \texthebrew{הצעת} prefix on otherwise identical strings
  \item Quote mark and whitespace variants
    (\texthebrew{על ידי} vs.\ \texthebrew{על-ידי})
\end{itemize}

\subsection*{Fuzzy Clustering Algorithm}

We applied fuzzy string matching (\texttt{rapidfuzz} token ratio, threshold = 88/100)
to cluster near-duplicates and assign a canonical representative per cluster
(the longest string in each cluster --- typically the most complete form).

The algorithm identified \textbf{23 non-trivial clusters}. After canonicalization,
the vocabulary reduced from 180 raw strings to \textbf{151 canonical topics}.

\subsection*{Manual Corrections}

At threshold 88, fuzzy matching produced \textbf{14 false-positive merges} --- cases where
two genuinely distinct topics scored above the similarity threshold due to shared legal
boilerplate phrasing. These were manually identified and corrected:

\begin{center}
\begin{tabular}{p{8cm}p{5cm}}
\toprule
\textbf{Incorrectly merged pair} & \textbf{Reason for split} \\
\midrule
Budget amendments 2022 \& 2023 & Different fiscal years \\
State Budget Law 2023 \& 2024 & Different fiscal years \\
Chapter \texthebrew{א'} (Preamble) \& Chapter \texthebrew{י'} \& Chapter \texthebrew{י"א} (Commencement) of the Economic Efficiency Law & Different chapters of same law \\
Chapter \texthebrew{א'} \& Chapter \texthebrew{כ"ו} of the Economic Plan Law & Different chapters of same law \\
\texthebrew{צו נכי המלחמה בנאצים} \& \texthebrew{צו נכי רדיפות הנאצים} & Different beneficiary populations \\
Tax ordinance amendment \#22 \& temporary order \#5 & Different legislative instruments \\
Ministers' salaries \& President's salary & Different officeholders \\
Public institution approval lists 2022 \& 2023 (section 61) & Different annual lists \\
Public institution approval lists 2022 \& 2023 (section 46) & Different annual lists \\
\texthebrew{פרק ט'} ``except section 36'' \& ``only section 36'' & Opposite inclusion scope \\
Temporary committee chair election \& standing committee chair election & Different committees \\
\bottomrule
\end{tabular}
\end{center}

% -----------------------------------------------------------------------
\section{Final Statistics}

\subsection*{Topic Vocabulary Summary}

\begin{center}
\begin{tabular}{lr}
\toprule
\textbf{Metric} & \textbf{Value} \\
\midrule
Total protocols & 213 \\
Raw unique topic strings & 180 \\
Genuine typo clusters merged & 23 \\
False-positive merges corrected & 14 \\
\textbf{Final canonical topics} & \textbf{151} \\
\midrule
Single-topic protocols & 191 \quad (89.7\%) \\
Multi-topic protocols & 22 \quad (10.3\%) \\
\bottomrule
\end{tabular}
\end{center}

\subsection*{Topic Frequency Distribution}

\begin{center}
\begin{tabular}{rrl}
\toprule
\textbf{Protocols} & \textbf{Topics} & \textbf{Share of topics} \\
\midrule
1  & 121 & 80.1\% \\
2  & 15  & 9.9\%  \\
3  & 4   & 2.6\%  \\
4  & 3   & 2.0\%  \\
6  & 2   & 1.3\%  \\
7  & 2   & 1.3\%  \\
9  & 1   & 0.7\%  \\
11 & 1   & 0.7\%  \\
12 & 1   & 0.7\%  \\
15 & 1   & 0.7\%  \\
\midrule
\textbf{Total} & \textbf{151} & \\
\bottomrule
\end{tabular}
\end{center}

\textbf{80\% of topics appear in only one protocol} (one-off discussions).
The remaining \textbf{30 recurring topics} (2+ protocols) are the core challenge
for streaming subject linking.

\subsection*{Most Recurring Topics}

\begin{center}
\begin{tabular}{rl}
\toprule
\textbf{\#} & \textbf{Topic} \\
\midrule
15 & \texthebrew{פרק ד' (מיסוי בשוק הדיור) --- חוק ההתייעלות הכלכלית} \\
12 & \texthebrew{שינויים בתקציב לשנת 2023} \\
11 & \texthebrew{פרק ט' (מיסים) --- חוק ההתייעלות הכלכלית} \\
9  & \texthebrew{פרק כ' (רשויות מקומיות, סעיפים 78--79) --- חוק ההתייעלות הכלכלית} \\
7  & \texthebrew{שינויים בתקציב לשנת 2022} \\
7  & \texthebrew{הצעת חוק התקציב לשנת הכספים 2023, מ/1610} \\
6  & \texthebrew{פרק כ' (רשויות מקומיות, סעיף 79)} \\
6  & \texthebrew{הצעת חוק לעידוד תעשייה עתירת ידע, מ/1577} \\
\bottomrule
\end{tabular}
\end{center}

% -----------------------------------------------------------------------
\section{Implications for Downstream Tasks}

\begin{itemize}
  \item \textbf{Topic \& segment extraction:} The \texttt{<{}< נושא >{}>} markers provide
    near-perfect segment boundary annotations in 96\% of protocols. These serve as
    silver labels and a strong rule-based baseline. The 4 marker-free files require
    model-based segmentation and are natural test cases.

  \item \textbf{Streaming subject linking evaluation:} Single-occurrence topics are
    trivially classified as ``new''. Evaluation signal comes entirely from the
    \textbf{30 recurring topics}. We recommend sampling the held-out evaluation set
    exclusively from these, ensuring coverage of both the high-frequency
    (budget law chapters) and low-frequency (2--4 occurrences) recurring cases.

  \item \textbf{Multi-topic protocols (22 files):} These sessions cover 2--3 distinct
    legislative matters. Correct segmentation is a strict prerequisite before linking ---
    conflating multiple topics into one embedding degrades recall on the linking task.
\end{itemize}

% -----------------------------------------------------------------------
\vspace{2em}
\hrule
\vspace{0.5em}
{\small
\textbf{Output files:}
\texttt{outputs/topics\_raw.json},
\texttt{outputs/topics\_canonical.json},
\texttt{outputs/topics\_clusters.json},
\texttt{outputs/topics\_corrections.json}.
\textbf{Code:} \texttt{src/preprocess\_topics.py}
}

\end{document}
